\section{Algorithm}
\label{sec:algorithm}

\subsection{High-Level Structure}

Algorithm~\ref{alg:nsga2} gives the top-level \nsga{} procedure for
redistricting.
The algorithm maintains a population of $N_{\mathrm{pop}} = 100$ valid plans
and evolves them over $N_{\mathrm{gen}} = 200$ generations.
Each generation applies non-dominated sort, crowding distance, tournament
selection, crossover, and mutation; the combined parent-plus-offspring pool is
truncated to $N_{\mathrm{pop}}$ by the same rank-then-crowding criterion.
After all generations, the Pareto front $\mathcal{F}_0$ of the final population
is returned.

\begin{algorithm}[t]
\caption{\nsga{}-\textsc{Redistrict}($G$, $\mathbf{p}$, $k$,
  $N_{\mathrm{pop}}$, $N_{\mathrm{gen}}$, $b$)}
\label{alg:nsga2}
\begin{algorithmic}[1]
\Require adjacency graph $G=(V,E)$, population vector $\mathbf{p}$,
  $k$ districts, $N_{\mathrm{pop}}$ plans, $N_{\mathrm{gen}}$ generations,
  base seed $b$
\State \textbf{// Initialise: $N_{\mathrm{pop}}$ METIS plans with distinct seeds}
\For{$i \gets 0$ \textbf{to} $N_{\mathrm{pop}}-1$}
  \State $P[i] \gets \textsc{MetisPlan}(G, \mathbf{p}, k,\; \textsc{InitSeed}(i, b))$
\EndFor
\For{$\mathrm{gen} \gets 0$ \textbf{to} $N_{\mathrm{gen}}-1$}
  \State $\mathbf{obj} \gets [\textsc{Evaluate}(P[i]) \;\forall\, i]$
    \Comment{$(\EC, \Dseats, \VRA)$ for each plan}
  \State $\mathbf{rank}, \mathbf{fronts} \gets \textsc{FastNonDominatedSort}(\mathbf{obj})$
  \State $\mathbf{crowd} \gets [\textsc{CrowdingDist}(\mathbf{fronts}[j], \mathbf{obj}) \;\forall\, j]$
  \State $\mathbf{parents} \gets \textsc{TournamentSelect}(P, \mathbf{rank}, \mathbf{crowd})$
  \State \textbf{// Crossover and mutation}
  \State $Q \gets [\,]$
  \For{$i \gets 0$ \textbf{to} $N_{\mathrm{pop}}/2 - 1$}
    \State $c \gets \textsc{Crossover}(\mathbf{parents}[2i], \mathbf{parents}[2i+1],\;
                                       G, \mathbf{p}, k,\; \textsc{CrossSeed}(\mathrm{gen}, i, b))$
    \State $c \gets \textsc{Mutate}(c,\; G, \mathbf{p},\; \varepsilon_{\mathrm{bal}},\;
                                    \textsc{MutSeed}(\mathrm{gen}, i, b))$
    \State $Q.\textsc{push}(c)$
  \EndFor
  \State \textbf{// Combine parent + offspring, truncate to $N_{\mathrm{pop}}$}
  \State $R \gets P \cup Q$; \quad $\mathbf{obj}_R \gets [\textsc{Evaluate}(R[i]) \;\forall\, i]$
  \State $\mathbf{rank}_R, \mathbf{fronts}_R \gets \textsc{FastNonDominatedSort}(\mathbf{obj}_R)$
  \State $\mathbf{crowd}_R \gets [\textsc{CrowdingDist}(\mathbf{fronts}_R[j], \mathbf{obj}_R) \;\forall\, j]$
  \State $P \gets \textsc{SelectTopN}(R, N_{\mathrm{pop}}, \mathbf{rank}_R, \mathbf{crowd}_R)$
\EndFor
\State \textbf{// Return Pareto front of final population}
\State $\mathbf{obj}_f \gets [\textsc{Evaluate}(P[i]) \;\forall\, i]$
\State $\mathbf{rank}_f, \mathbf{fronts}_f \gets \textsc{FastNonDominatedSort}(\mathbf{obj}_f)$
\State \Return \textsc{ParetoResult} with $\{P[i] : \mathbf{rank}_f[i] = 0\}$
\end{algorithmic}
\end{algorithm}

\subsection{Non-Dominated Sort: \textsc{FastNonDominatedSort}}

The fast non-dominated sort assigns a Pareto rank to each plan in
$O(M N^2)$ time, where $M = 3$ is the number of objectives and
$N = |P|$ is the population size.
For $N = 200$ (combined parent-plus-offspring pool), this is
$3 \times 40{,}000 = 120{,}000$ comparisons per generation --- negligible
compared to the cost of evaluating plans.

Algorithm~\ref{alg:fnds} gives the full procedure.

\begin{algorithm}[t]
\caption{\textsc{FastNonDominatedSort}($\mathbf{obj}$)}
\label{alg:fnds}
\begin{algorithmic}[1]
\State $n \gets |\mathbf{obj}|$; \quad
  $\mathbf{cnt} \gets [0]^n$; \quad
  $S \gets [\emptyset]^n$
  \Comment{$\mathbf{cnt}[p]$ = how many plans dominate $p$}
\For{$p \gets 0$ \textbf{to} $n-1$}
  \For{$q \gets 0$ \textbf{to} $n-1$}
    \If{$\textsc{Dominates}(\mathbf{obj}[p], \mathbf{obj}[q])$}
      $S[p] \gets S[p] \cup \{q\}$
    \ElsIf{$\textsc{Dominates}(\mathbf{obj}[q], \mathbf{obj}[p])$}
      $\mathbf{cnt}[p] \mathrel{+}= 1$
    \EndIf
  \EndFor
\EndFor
\State $\mathbf{fronts}[0] \gets \{p : \mathbf{cnt}[p] = 0\}$; \quad $i \gets 0$
\While{$\mathbf{fronts}[i] \neq \emptyset$}
  \State $\mathbf{fronts}[i+1] \gets \emptyset$
  \For{$p \in \mathbf{fronts}[i]$}
    \For{$q \in S[p]$}
      \State $\mathbf{cnt}[q] \mathrel{-}= 1$
      \If{$\mathbf{cnt}[q] = 0$}
        $\mathbf{fronts}[i+1] \gets \mathbf{fronts}[i+1] \cup \{q\}$
      \EndIf
    \EndFor
  \EndFor
  \State $i \mathrel{+}= 1$
\EndWhile
\State \Return $\mathbf{rank}$ where $\mathbf{rank}[p] = i$ iff $p \in \mathbf{fronts}[i]$,
  and $\mathbf{fronts}$
\end{algorithmic}
\end{algorithm}

The dominance check \textsc{Dominates}$(A, B)$ returns \texttt{true} iff
$A[m] \leq B[m]$ for all $m \in \{\EC, \Dseats, \VRA\}$ and
$A[m] < B[m]$ for some $m$.
This is a coordinate-wise partial order; ties on all objectives are not
dominance (neither plan dominates the other).

\subsection{Crowding Distance}

Within each Pareto front, crowding distance measures isolation in objective
space.
Plans at the extremes of any objective (minimum or maximum value on that
objective within the front) receive infinite crowding distance and are
always preferred during selection.
Interior plans receive crowding distance proportional to the normalised gap
between their two neighbours on each objective.

\begin{algorithm}[t]
\caption{\textsc{CrowdingDist}($F$, $\mathbf{obj}$)}
\label{alg:crowd}
\begin{algorithmic}[1]
\State $n \gets |F|$; \quad $\mathbf{dist} \gets [0.0]^n$
\For{each objective $m \in \{\EC, \Dseats, \VRA\}$}
  \State $\sigma \gets \textsc{SortByObj}(F, m)$
    \Comment{indices of $F$ sorted by $\mathbf{obj}[\cdot][m]$ ascending}
  \State $\mathbf{dist}[\sigma[0]] \gets +\infty$; \quad
    $\mathbf{dist}[\sigma[n-1]] \gets +\infty$
  \State $r \gets \mathbf{obj}[\sigma[n-1]][m] - \mathbf{obj}[\sigma[0]][m]$
  \If{$r = 0$} \textbf{continue} \EndIf
    \Comment{all plans identical on objective $m$}
  \For{$i \gets 1$ \textbf{to} $n-2$}
    \State $\mathbf{dist}[\sigma[i]] \mathrel{+}=
      (\mathbf{obj}[\sigma[i+1]][m] - \mathbf{obj}[\sigma[i-1]][m]) / r$
  \EndFor
\EndFor
\State \Return $\mathbf{dist}$
\end{algorithmic}
\end{algorithm}

\subsection{Tournament Selection}

Binary tournament selection draws two plans at random and returns the winner
under a two-tier criterion: prefer lower Pareto rank; break ties by higher
crowding distance.

\begin{algorithm}[t]
\caption{\textsc{TournamentSelect}($P$, $\mathbf{rank}$, $\mathbf{crowd}$)}
\label{alg:tournament}
\begin{algorithmic}[1]
\State $\mathbf{parents} \gets [\,]$
\For{$\_ \gets 0$ \textbf{to} $N_{\mathrm{pop}}-1$}
  \State $(a, b) \gets \textsc{RandomPair}(\{0,\ldots,|P|-1\})$
  \If{$\mathbf{rank}[a] < \mathbf{rank}[b]$}
    $\mathbf{parents}.\textsc{push}(P[a])$
  \ElsIf{$\mathbf{rank}[b] < \mathbf{rank}[a]$}
    $\mathbf{parents}.\textsc{push}(P[b])$
  \ElsIf{$\mathbf{crowd}[a] > \mathbf{crowd}[b]$}
    $\mathbf{parents}.\textsc{push}(P[a])$
  \Else
    $\mathbf{parents}.\textsc{push}(P[b])$
  \EndIf
\EndFor
\State \Return $\mathbf{parents}$
\end{algorithmic}
\end{algorithm}

\textsc{SelectTopN} uses the same two-tier criterion iteratively: add entire
fronts until the budget $N_{\mathrm{pop}}$ would be exceeded, then fill the
remaining slots from the next front sorted by crowding distance descending.

\subsection{Crossover: ReCom-Style District Merge}

The crossover operator adapts the ReCom redistricting step~\citep{deford2021}
to the genetic algorithm setting.
It merges two adjacent districts from parent $A$, then redraws the merged
region using a spanning-tree proposal seeded from the crossover seed.
Unlike full ReCom, no Metropolis-Hastings acceptance ratio is needed:
NSGA-II's selection pressure replaces the Markov chain acceptance criterion.

\begin{algorithm}[t]
\caption{\textsc{Crossover}($A$, $B$, $G$, $\mathbf{p}$, $k$, $s$)}
\label{alg:crossover}
\begin{algorithmic}[1]
\State $(d_1, d_2) \gets \textsc{RandomAdjacentDistrictPair}(A, G, s)$
  \Comment{seed $s$ for reproducibility}
\State $R \gets \{t \in V : A[t] \in \{d_1, d_2\}\}$
  \Comment{merge districts $d_1$ and $d_2$}
\For{attempt $\gets 1$ \textbf{to} $5$}
  \State $T \gets \textsc{WilsonSpanningTree}(G[R],\; s \oplus \mathrm{attempt})$
    \Comment{random spanning tree of subgraph induced by $R$}
  \State $e^* \gets \textsc{FindBalancedCutEdge}(T, \mathbf{p}, \varepsilon_{\mathrm{bal}})$
    \Comment{cut $T$ into two components within balance tolerance}
  \If{$e^* \neq \texttt{None}$}
    \State $(R_1, R_2) \gets \textsc{Split}(T, e^*)$
    \State $C \gets A$; \quad $C[t] \gets d_1$ for $t \in R_1$; \quad
      $C[t] \gets d_2$ for $t \in R_2$
    \State \Return $C$
  \EndIf
\EndFor
\State \Return $A$
  \Comment{all retries failed: return parent $A$ unchanged}
\end{algorithmic}
\end{algorithm}

The fallback to parent $A$ (line 11) is intentional.
\nsga{} tolerates a fraction of no-op crossovers without degenerating,
provided the mutation operator is active.
If the crossover validity rate falls below 20\%, Phase~2 will investigate
improved proposals (e.g., preferring adjacent district pairs with nearly
balanced populations, which are more likely to yield a balanced spanning-tree
cut).

\subsection{Mutation: Single Boundary-Tract Flip}

The mutation operator makes exactly one boundary-tract flip, reusing the
Flip chain logic from the \texttt{bisect-cli} codebase.
The invariant is that all plans in the population are always valid: the
mutation either produces a valid flipped plan or returns the original
plan unchanged.

\begin{algorithm}[t]
\caption{\textsc{Mutate}($\pi$, $G$, $\mathbf{p}$, $\varepsilon_{\mathrm{bal}}$, $s$)}
\label{alg:mutate}
\begin{algorithmic}[1]
\State $B \gets \{t \in V : \exists\, u \in N(t) \text{ s.t.\ } \pi[u] \neq \pi[t]\}$
  \Comment{boundary tracts}
\State $t \gets \textsc{RandomElement}(B, s)$
\State $D_{\mathrm{adj}} \gets \{\pi[u] : u \in N(t),\; \pi[u] \neq \pi[t]\}$
  \Comment{adjacent target districts}
\State $d_{\mathrm{new}} \gets \textsc{RandomElement}(D_{\mathrm{adj}}, s)$
\State $\pi' \gets \pi$; \quad $\pi'[t] \gets d_{\mathrm{new}}$
\State \textbf{// Validate: population balance}
\If{$|\mathrm{pop}(d_{\mathrm{new}}, \pi') - P_{\mathrm{ideal}}| > \varepsilon_{\mathrm{bal}} \cdot P_{\mathrm{ideal}}$
    $\;\vee\;$
    $|\mathrm{pop}(\pi[t], \pi') - P_{\mathrm{ideal}}| > \varepsilon_{\mathrm{bal}} \cdot P_{\mathrm{ideal}}$}
  \State \Return $\pi$
    \Comment{balance violated: return unchanged}
\EndIf
\State \textbf{// Validate: contiguity of the donor district}
\If{\textbf{not} $\textsc{IsConnected}(G[\{u : \pi'[u] = \pi[t]\}])$}
  \State \Return $\pi$
    \Comment{donor district disconnected: return unchanged}
\EndIf
\State \Return $\pi'$
\end{algorithmic}
\end{algorithm}

The mutation makes exactly one attempt.
If the selected flip is invalid (fails balance or contiguity check), the
plan is returned unchanged.
This single-attempt design is intentional: multiple retries would bias
the effective mutation rate toward boundary tracts with valid flips,
introducing a selection pressure not accounted for by the NSGA-II framework.

\subsection{Seeding Specification}

All stochastic operations are derived from \texttt{base\_seed} via SHA-256
with domain-separated prefixes.
Full reproducibility is guaranteed given only \texttt{base\_seed}.

\begin{definition}[Seed functions]
\label{def:seeds}
Let $\|\cdot\|$ denote byte concatenation and let $\mathbf{SHA256}(x)[0:8]$
denote the first 8 bytes of the SHA-256 hash of $x$, interpreted as a
little-endian \texttt{u64}.

\begin{align*}
\textsc{InitSeed}(i,\, b) &= \mathbf{SHA256}(
  \texttt{"PARETO\_INIT\_"} \,\|\, i_{\mathrm{u32le}} \,\|\, \texttt{"\_"}
  \,\|\, b_{\mathrm{u64le}})[0:8], \\
\textsc{CrossSeed}(g,\, i,\, b) &= \mathbf{SHA256}(
  \texttt{"PARETO\_CROSS\_"} \,\|\, g_{\mathrm{u32le}} \,\|\, \texttt{"\_"}
  \,\|\, i_{\mathrm{u32le}} \,\|\, \texttt{"\_"} \,\|\, b_{\mathrm{u64le}})[0:8], \\
\textsc{MutSeed}(g,\, i,\, b) &= \mathbf{SHA256}(
  \texttt{"PARETO\_MUT\_"} \,\|\, g_{\mathrm{u32le}} \,\|\, \texttt{"\_"}
  \,\|\, i_{\mathrm{u32le}} \,\|\, \texttt{"\_"} \,\|\, b_{\mathrm{u64le}})[0:8].
\end{align*}
\end{definition}

The prefix strings are version-locked: any change to the proposal algorithm
must change the prefix to prevent silent seed compatibility across versions.
The three prefixes differ in byte length (\texttt{PARETO\_INIT\_} = 12 bytes,
\texttt{PARETO\_CROSS\_} = 13 bytes, \texttt{PARETO\_MUT\_} = 11 bytes) and
in content, ensuring
$\textsc{InitSeed}(i, b) \neq \textsc{CrossSeed}(g, i, b) \neq \textsc{MutSeed}(g, i, b)$
for all $i$, $g$, $b$ --- no seed collisions between the three operations.

\subsection{SMC-Pareto Alternative}

As an alternative to running the full \nsga{} algorithm, practitioners who
have already produced an SMC ensemble~\citep{dellaLibera2026smc} can project
the particle cloud onto objective space and extract the Pareto front.

\begin{verbatim}
bisect pareto \
  --method smc-project \
  --smc-input nc_smc_2020.ndjson \
  --output nc_pareto_smc_2020.ndjson
\end{verbatim}

The \texttt{smc\_project} module reads the SMC NDJSON file, evaluates all
three objectives for each particle (weighted by SMC importance weight),
runs \textsc{FastNonDominatedSort} on the particle population, and outputs
the Pareto front in the same NDJSON format as \nsga{}.

\paragraph{Comparison with NSGA-II.}
SMC-Pareto is substantially faster (seconds to minutes) because it reuses
an existing ensemble.
However, SMC samples from the near-uniform distribution over valid plans
weighted by compactness; it may undersample extreme regions of the Pareto
front --- plans with very high compactness at the cost of poor VRA, or plans
with perfect VRA at the cost of poor compactness.
\nsga{} actively searches for Pareto-optimal diversity and is expected to
find plans in these extreme regions.
Phase~2 will measure whether \nsga{} finds Pareto regions that the SMC
projection misses.

\subsection{CLI Interface}

The \nsga{} redistricting algorithm is exposed as a top-level
\texttt{bisect pareto} subcommand:

\begin{verbatim}
# Run NSGA-II for NC congressional (k=14)
bisect pareto \
  --state NC --year 2020 \
  --population 100 \
  --generations 200 \
  --base-seed 42 \
  --output nc_pareto_2020.ndjson

# Enacted plan dominance test
bisect pareto \
  --state NC --year 2020 \
  --population 100 \
  --generations 200 \
  --base-seed 42 \
  --enacted-plan configs/enacted_nc.yml \
  --output nc_pareto_2020.ndjson
\end{verbatim}

Output is NDJSON: one line per Pareto-optimal plan, followed by an optional
enacted-plan record, followed by a metadata record containing
$N_{\mathrm{pop}}$, $N_{\mathrm{gen}}$, Pareto front size, \texttt{base\_seed},
runtime, and \texttt{file\_sha256} (SHA-256 of all preceding lines).
The \texttt{file\_sha256} field provides a tamper-evident audit chain over
the returned frontier.

Table~\ref{tab:scaling} gives expected parameter scaling and runtime.

\begin{table}[t]
\centering
\caption{%
  NSGA-II parameter scaling by state size.
  Runtimes are approximate and hardware-dependent; all timing estimates
  assume a single sequential thread on a 2025-era 8-core workstation.
  Pareto front size grows with $k$ because the objective space becomes larger.
}
\label{tab:scaling}
\renewcommand{\arraystretch}{1.3}
\begin{tabular}{@{}lrrrrl@{}}
\toprule
State size & $N_{\mathrm{pop}}$ & $N_{\mathrm{gen}}$ &
  Runtime (sequential) & Pareto front size \\
\midrule
Small ($k \leq 5$, e.g.\ VT) & 50 & 100 & $<5$ min & 3--8 \\
Medium ($k = 8$--14, e.g.\ WI, NC) & 100 & 200 & 15--60 min & 10--30 \\
Large ($k \geq 30$, e.g.\ TX, CA) & 200 & 500 & 4--12 hours & 20--60 \\
\bottomrule
\end{tabular}
\end{table}
